\subsection{Every mixed line with primitive projections}

\begin{theorem}[DP42: sharp primitive mixed-line slopes]\label{DP42}
Assume $p$ is odd and $h\ge3$, and use exactly DP41's group, dual, coordinates and invariant primitive frame over $\mathcal O=W(\F_q)$. Let $u\in\mathcal O^\times$ and let $\mathbf v=(v_1,\ldots,v_d)\in\mathcal O^d$ be primitive. Put
\[
 j=\min\{i:v_i\in\mathcal O^\times\},\qquad k=h-j,\qquad
 N_k=p^k+1,\qquad S_k=\frac{p^k-1}{p-1},
\]
and
\[
 w_k=\frac{p^k}{(p-1)(p^k+1)},\qquad
 c_k^{\mathrm{path}}=\frac1{(p-1)(p^k+1)}.
\]
For an integral $U\in\mathcal O[[S,\mathbf Y]]$ assume only that
\[
 V(t)=U(E_0(ut),E_H(\mathbf vt))
\]
is not identically zero. It may vanish at the origin. For $\lambda\in\mathcal O^\times$ set
$\Theta(t)=V(t)\exp(\lambda\beta(ut,\mathbf vt))=\sum a_nt^n$.
For every rational $w\ge0$,
\[
 \nu_w(\Theta):=\inf_n(\vp(a_n)+wn)=
 \begin{cases}
 -\infty,&w<w_k,\\
 \nu_w(V)\in[0,\infty),&w\ge w_k.
 \end{cases}
\]
The finite value is zero if $U(0)$ is a unit. In particular
\begin{equation}\label{eq:pathsharp}
 \vp(n!a_n)\ge c_k^{\mathrm{path}}n-\frac{s_p(n)}{p-1},
 \qquad
 \liminf_{n\to\infty}\frac{\vp(n!a_n)}n=c_k^{\mathrm{path}}.
\end{equation}
The logarithmic restriction also has the exact congruence
\begin{equation}\label{eq:pathcoefficient}
 p^{S_k}[t^{N_k}]\beta(ut,\mathbf vt)
       \equiv(-1)^{k-1}u v_j^{p^k}\pmod p.
\end{equation}
The integer $k$ is the restricted-$p$ nilpotency length of the reduced dual direction. Thus the largest slope among all these primitive mixed lines is
$c_1^{\mathrm{path}}=1/(p^2-1)$, independent of $h$; longer restricted-$p$ chains give strictly smaller slopes.
\end{theorem}

\begin{proof}
Put $s=1/(p-1)$. We isolate the full weighted initial form, including its behavior under inverse logarithms; a first pole alone would not determine the exponential radius.
For rational
\[
 w\in I_k:=\left[\frac{S_k}{p^k+1},\,w_k\right],
 \qquad \delta=s-w,
\]
define
\[
 \omega_0=w,\qquad \omega_i=s-p^{i-j}\delta\quad(1\le i\le d).
\]
Then $\omega_j=w$,
\[
 \frac{s}{p^k+1}\le\delta\le\frac{2s}{p^k+1}<s,\qquad
 0\le\tau:=2s-(p^k+1)\delta=(p^k+1)w-S_k\le s.
\]
All weights are positive and exceed DP41's original contraction weights:
since $p\ge3$,
\[
 \delta\le\frac{2s}{p^k+1}\le\frac1{p^k+1}
 <\frac1{p^k-p^{-j}}=p^j\rho,
 \qquad \omega_i-r_i=p^i(\rho-p^{-j}\delta)>0.
\]
Also $w\ge1/(p+1)>\rho$. These imply
\begin{equation}\label{eq:pathstrict}
 w(q-1)>1,\qquad p^{h-i}\omega_i(q-1)>d.
\end{equation}
For the second inequality it suffices to replace $\omega_i$ by $r_i$:
the resulting expression is at least
$(q-p)/(p-1)=p+\cdots+p^{h-1}>d$.

\emph{Integral logarithms and their inverse reductions.}
Choose $\pi^e=p$ clearing all weight denominators and put
$\alpha_i=\pi^{e\omega_i}$, $\gamma=\alpha_0=\alpha_j$.
The contracted scalar logarithm reduces to the identity, since each nonlinear coefficient has valuation $(q^a-1)w-a>0$.
For the dual arrows,
\[
 p\omega_{i-1}-\omega_i=1\quad(i\ge2),\qquad
 p^2\omega_d-\omega_1
       =p+1-p^{1-j}(q-1)\delta>1.
\]
A composition of at least two arrows has weight increment at least
$1+p+\cdots+p^{a-1}>a$, so all higher iterates divided by $p^a$ vanish in reduction. The two contracted logarithms are integral with identity differential and integral inverses, and the dual reduction is exactly
\[
 \overline{\widetilde L}_1=Y_1,\qquad
 \overline{\widetilde L}_i=Y_i+Y_{i-1}^p\quad(i\ge2).
\]
The scaled additive input in coordinate $i$ is
$z_i(T)=v_i\gamma\alpha_i^{-1}T$.
For $i>j$ its valuation is positive because $\omega_i<w$.
For $i<j$ we have $v_i\in p\mathcal O$, whereas
$0<\omega_i-w=\delta(1-p^{i-j})<s<1$, so its valuation is again positive. The unramified integral lattice is used here. Only $z_j=v_jT$ survives. Successive inversion therefore gives
\begin{equation}\label{eq:pathinversereduction}
 \overline{\widetilde E}_i(z(T))=
 \begin{cases}
 0,&i<j,\\
 (-1)^{i-j}\bar v_j^{p^{i-j}}T^{p^{i-j}},&i\ge j,
 \end{cases}
 \qquad
 \overline{\widetilde E}_0(uT)=\bar uT.
\end{equation}
Oddness of $p$ preserves these signs under $p$th powers.

\emph{The unique initial term of the primitive logarithm.}
Apply the weights $\boldsymbol\omega$ to \eqref{eq:fullprimitive}, and put
$\mu=\alpha_0\alpha_d^p/p$, of valuation $\tau$.
For its first family set
\[
 A_i=p^iw+\omega_i=s+p^i\bigl(s-\delta(1+p^{-j})\bigr).
\]
The bracket is positive since $p^k+1>2(1+p^{-j})$; thus $A_i$ increases with $i$. Its smallest baseline has excess
\[
 A_1-1-\tau
       =\delta\bigl(p^k-p+1-p^{1-j}\bigr)>0.
\]
For $k=1$ the bracket is $1-p^{1-j}>0$ because $h>2$; for $k>1$ it is positive even at $j=1$. A general first-family term has excess over its baseline
\[
 (p^n-1)A_i-n+p^{n+i}w(q^l-1)-l.
\]
Both parts are nonnegative, strictly positive for a positive corresponding index, using $A_i>1$ and \eqref{eq:pathstrict}. Thus the first family lies strictly above $\tau$.

In the product family put $e_i=h-i$. Its baseline excess is
\[
 w+p^{e_i}\omega_i-e_i-\tau
       =\frac{p^{e_i}-1}{p-1}-e_i.
\]
This is zero precisely for $i=d$, and positive otherwise. Its further excess at indices $a,l$ is
\[
 w(q^a-1)-a+p^{e_i}\omega_i(q^l-1)-dl,
\]
strictly positive if either index is positive, by \eqref{eq:pathstrict}. Consequently the unique surviving term is $SY_d^p/p$, and
\begin{equation}\label{eq:pathinitial}
 \overline{\mu^{-1}b(\alpha_0S,\alpha_1Y_1,\ldots,\alpha_dY_d)}
       =SY_d^p.
\end{equation}
No cancellation can occur at this weight.

Composition of \eqref{eq:pathinitial} with the integral inverse maps commutes with reduction. Equation \eqref{eq:pathinversereduction} now gives the entire restricted initial form
\[
 \overline{\mu^{-1}\beta(u\gamma T,\mathbf v\gamma T)}
       =(-1)^{k-1}\bar u\bar v_j^{p^k}T^{p^k+1}.
\]
Its derivative is nonzero in characteristic $p$. With the common chosen powers of $\pi$ one has exactly
$\mu=\gamma^{p^k+1}/p^{S_k}$.
Comparing that coefficient proves \eqref{eq:pathcoefficient}: its original coefficients lie in the unramified field, so positive valuation of the difference implies divisibility by $p$.

\emph{Exponential threshold and the prefactor.}
For $w\in I_k$ below $w_k$, the logarithmic coefficient infimum is
$\tau\in[0,s)$ with the displayed nonzero-derivative initial form.
DP41's fraction-field multiplier lemma therefore shows that every nonzero bounded factor leaves $\exp(\lambda\beta(u\gamma T,\mathbf v\gamma T))$ unbounded.
The factor $V(\gamma T)$ is integral by the integral contracted inputs and inverse maps, and is nonzero by hypothesis. Thus $\nu_w(\Theta)=-\infty$ on this interval; monotonicity gives the same conclusion for every smaller nonnegative $w$.

At $w=w_k$ the logarithm has coefficient infimum $s$, so its exponential is integral with coefficient infimum zero. Gauss multiplicativity gives
$\nu_{w_k}(\Theta)=\nu_{w_k}(V)\in[0,\infty)$.
For larger $w$ both factors remain integral by monotonicity and the exponential retains constant one, proving the asserted finite equality. A unit $U(0)$ makes this infimum zero.
In particular $\vp(a_n)\ge-w_kn$. If $\liminf\vp(a_n)/n$ were greater than $-w_k$, some rational $w<w_k$ would make the weighted infimum finite, a contradiction. Therefore $\liminf\vp(a_n)/n=-w_k$.
Legendre's formula and $s_p(n)/n\to0$ give \eqref{eq:pathsharp}. Zero coefficients have valuation $+\infty$; unboundedness already forces infinitely many nonzero ones.

\emph{Restricted-$p$ interpretation.}
Along dual axis $i<d$, the logarithm through degree $p$ is
$te_i+t^pe_{i+1}/p+O(t^{p+1})$.
The translation coefficient is
$D_i^p/p!+D_{i+1}/p$, an integral operator. Thus modulo $p$,
$\bar D_i^{[p]}$ is a nonzero scalar multiple of $\bar D_{i+1}$.
Along the last axis the first nonlinear logarithm term has degree $p^2$, so $D_d^p/p!$ is integral and $\bar D_d^{[p]}=0$.
Invariant directions commute, and their restricted operation is additive and Frobenius-semilinear. The first nonzero coordinate $j$ of the reduced vector therefore survives exactly $k-1=d-j$ iterations and dies on the next. Later coordinates cannot cancel its leading surviving component. This proves the intrinsic interpretation and, since $s/(p^k+1)\le s/(p+1)=1/(p^2-1)$, the maximum assertion.
\end{proof}

The result covers every mixed logarithmic line whose two projections are primitive in these unramified lattices, with any integral prefactor having nonzero restriction. It strengthens DP36's sufficient last-axis residual $\rho_h$ to the exact $\rho_2$. It does not cover a zero or nonprimitive projection, a nonunit exponential parameter, another prescribed gauge, or arbitrary selected harmonic differential polynomials. In particular it is a local obstruction to a proposed improvement on these paths, not a nonexistence theorem for every shorter global correspondence.

